\chapter[Planteamiento]{Planteamiento y decisión sobre las soluciones alternativas}

La etapa de diseño del sistema DECOR requirió evaluar varias alternativas arquitectónicas antes de fijar la solución definitiva. En este capítulo se describen las opciones consideradas para cada componente del sistema y se justifica la elección adoptada.

\section[Alternativas para la extrapolación]{Alternativas para la extrapolación de la reverberación tardía}

La cola de una RIR comparte tres características físicas universales: decaimiento exponencial, coloración espectral dependiente de los materiales y un componente estocástico de alta frecuencia. Esto abre dos familias de enfoques:

\begin{itemize}
	\item \textbf{Regresión directa de la señal temporal:} El modelo predice directamente las muestras de audio. Permite el uso de cualquier función de pérdida temporal o espectral, pero supone un objetivo de alta dimensionalidad (48000 muestras) con elevada varianza aleatoria, haciendo el aprendizaje inestable.
	\item \textbf{Modelo físico paramétrico:} El modelo predice parámetros de un modelo analítico de decaimiento (p. ej., coeficientes de decaimiento exponencial por banda de octava) que luego sintetizan la señal. Reduce el espacio de salida a unas pocas decenas de números, garantizando plausibilidad física por construcción.
\end{itemize}

Se optó por el \textbf{modelo físico paramétrico} por su estabilidad durante el entrenamiento y su interpretabilidad, siguiendo el marco propuesto en \cite{muhammad2025}.

\section[Alternativas para el codificador]{Alternativas para el codificador de la cabeza}

Para extraer el vector latente a partir de los primeros 50~ms de la RIR se consideraron tres arquitecturas:

\begin{table}[H]
	\centering
	\caption{Comparativa de arquitecturas de codificador consideradas}
	\begin{tabular}{|p{3cm}|p{5cm}|p{5cm}|}
		\hline
		\textbf{Arquitectura} & \textbf{Ventajas} & \textbf{Inconvenientes} \\ \hline
		LSTM bidireccional & Modela dependencias temporales nativas & Secuencial, difícil de paralelizar; costoso en señales largas \\ \hline
		Transformer & Atención global, altamente paralelizable & Coste cuadrático en longitud de secuencia sin optimizaciones \\ \hline
		CNN residual (elegida) & Paralelizable, campo receptivo controlado, estable con BN & Requiere un número elevado de bloques para cubrir toda la cabeza \\ \hline
	\end{tabular}
\end{table}

Se eligió la \textbf{CNN residual con stride} (\texttt{DecorEncoder}) porque combina campo receptivo suficientemente amplio ($>100.000$ muestras con 14 bloques de stride 2), paralelización completa en GPU y estabilidad de entrenamiento gracias a \textit{BatchNorm} y conexiones residuales.

\section{Alternativas para la función de pérdida}

La pérdida es un factor crítico para la calidad perceptual. Las alternativas estudiadas fueron:

\begin{itemize}
	\item \textbf{MSE en tiempo:} Simple pero ciega a la percepción espectral; penaliza errores de fase irrelevantes perceptualmente.
	\item \textbf{MAE en EDC:} Favorece la corrección de la envolvente energética; no penaliza errores espectrales de coloración.
	\item \textbf{MSTFT (elegida):} Combina pérdida de convergencia espectral y pérdida de log-magnitud en cuatro resoluciones temporales/frecuenciales, siendo la más alineada con la percepción humana y con el estado del arte \cite{Lin25}.
\end{itemize}

La implementación final (\texttt{models/loss.py}) utiliza resoluciones FFT de $\{64, 512, 2048, 8192\}$ con saltos $\{32, 256, 1024, 4096\}$, cubriendo desde detalles de ataque hasta la macro-envolvente del decaimiento.

\section{Decisión sobre el dataset de entrenamiento}

Se evaluaron dos fuentes de datos:

\begin{enumerate}
	\item \textbf{Dataset sintético propio:} 6000 salas rectangulares generadas con Pyroomacoustics (ISM). Ventaja: configuración controlada, trazabilidad total. Inconveniente: distribución acotada a geometrías simples.
	\item \textbf{BIRD \cite{Grondin25}:} Dataset real con 89 splits y miles de RIR medidas en recintos reales de todo tipo. Mayor diversidad geométrica y de materiales.
\end{enumerate}

Se adoptó \textbf{BIRD como fuente principal} de entrenamiento por su diversidad, reservando el dataset sintético para estudios de ablación y validación de la pipeline de datos.